% !TeX root = ../main.tex

\chapter{Introduction}\label{chap:introduction}

\section{Background}

Open with the problem your field cares about, in language a first-year graduate
student in a neighbouring department could follow. Replace this paragraph.

\lipsum[1][1-6]

\section{Motivation}

State the gap: what is unknown, unmeasured, or unresolved, and why it matters
that somebody close it. A committee reads this section to decide whether the
rest of the thesis was worth writing.

\lipsum[2][1-5]

\section{Contributions}

List what this thesis adds. Keep the claims narrow enough that
Chapter~\ref{chap:results} can actually support them.

\begin{itemize}
  \item The first contribution.
  \item The second contribution.
  \item The third contribution.
\end{itemize}

\section{Thesis Organization}

% \ref{...} 會被替換成章號。編譯兩次才會正確：第一次把章號寫進 main.aux，
% 第二次才讀得到。latexmk 會自動處理。
% \ref{...} expands to the chapter number. It takes two compiler passes to
% resolve — the first writes main.aux, the second reads it. latexmk handles
% this automatically, which is why the build uses it.
Chapter~\ref{chap:template} shows what this template can typeset and how to
drive it. Chapter~\ref{chap:methods} describes the methods.
Chapter~\ref{chap:results} reports the results, and
Chapter~\ref{chap:discussion} discusses what they mean, where they fall short,
and what should come next.
